\section{The PDF layer}
\label{sec:pdf-layer}

Five operations touch bytes: parse, incremental write (annotations), whole-document rewrite
(page operations), merge, and flatten (annotations baked into page content). They share a
parser and differ in what they are allowed to assume.

\subsection{Parsing}

The parser (\texttt{src/pdf/syntax.ts}, \texttt{document.ts}) handles the parts of ISO 32000
that real documents use: classic cross-reference tables, cross-reference streams, object
streams, incremental update chains via \texttt{/Prev}, and the standard filters for the
streams it must decode (none are decoded for a verbatim copy). It retains for every object
both the parsed value and, when the object was stored uncompressed, the exact source text.

The parsed value distinguishes streams by carrying their raw bytes:
\begin{lstlisting}[language={}]
type PdfStreamValue = {
  kind: 'stream';
  dict: PdfDict;
  raw: Uint8Array;   // exactly as it appears in the file, before any filter
};
\end{lstlisting}
This one field is what makes merging possible without a codec: a stream from another document
can be written back with its bytes untouched, because the only thing that needed rewriting
was the dictionary's references. See \S\ref{sec:merge}.

Two parser details caused defects worth recording, because both are invisible until a file is
written and re-read. First, an end-of-stream marker has two positions --- where the
\texttt{endstream} keyword begins and where its preceding newline ends --- and slicing content
with the trimmed one left a trailing newline that looked like the start of the next token.
Second, object numbers are not necessarily contiguous: a file may reference an object it does
not define, and the parser must distinguish ``missing'' from ``null'', since a page's
\texttt{/Annots} may legitimately be absent.

\subsection{The incremental annotation writer}
\label{sec:incremental}

Saving annotations appends objects to the existing file, rewrites the page dictionaries that
gained or lost annotations, and appends a cross-reference section whose trailer carries
\texttt{/Prev} pointing at the previous one. The original bytes are never modified, so the
operation is reversible by truncation and a reader that ignores the new section still sees a
valid document.

The invariant that governs everything is how a page's annotation array is computed. Let
$E$ be the object numbers currently listed in the page's \texttt{/Annots}, $R$ the object
numbers this writer produced for that page on the previous save, and $N$ the object numbers
of the annotations now in the model. Then

\begin{lstlisting}[language={}]
annots = [ refs(E) \ refs(R) ] ++ [ refs(N) ]
\end{lstlisting}

\noindent
and the page dictionary is rewritten whenever $N \cup R \neq \emptyset$ --- including when
$N = \emptyset$, because a page whose annotations were all deleted has no other reason to be
rewritten, and its \texttt{/Annots} would otherwise keep pointing at them.

\paragraph{What goes wrong without it.} The first implementation appended:
$\texttt{annots} = E \cup N$. Two symptoms followed, both reported by a user before any test
noticed.

\begin{enumerate}
  \item \textbf{Duplication.} Every save appended a fresh copy of every annotation on the page
        and listed both. Because the copies are byte-identical and coincide geometrically,
        the page looks correct; the file grows, and the reader may select one copy while
        deleting another. This reads to a user as ``the delete button does nothing''.
  \item \textbf{Resurrection.} A deleted annotation was removed from the model but its object
        was still listed in $E$, so the next save re-published it. The mark could not be
        removed from the document by any sequence of operations in the application.
\end{enumerate}

The fix records, per page, the object numbers the writer produced, keeps that map on the
document session, and passes it back on the next save as $R$; the writer uses the numbers it
produced this time as the new $R$. Objects superseded by earlier saves remain in the file as
orphans --- unreferenced, legal, and the price of never rewriting the original bytes.
\S\ref{sec:verification} describes the regression test, and a note on the test itself: the
first version of that test counted annotations in the file by calling the writer, which is a
write, so the counting added copies. It now fetches the saved file and resolves the page's
\texttt{/Annots}, counting what the file points at rather than what it contains.

\paragraph{Interoperability versus fidelity.} Every annotation is written twice, in a sense:
once as a standard annotation object that other readers display, and once as a private
\texttt{/PLP} dictionary holding this application's exact kind, stroke width, dash pattern and
font size. Standard subtypes cannot express all of it --- a text box and a callout are both
\texttt{/FreeText}, a replaced word is a \texttt{/StrikeOut} --- so the private entry is what
makes a round trip lossless. Section~\ref{sec:annot-model} discusses the cost of this duality.

\subsection{Whole-document rewrite}
\label{sec:rewrite}

Delete, insert, reorder, extract, crop and rotate pages cannot be expressed as an appending
update if they move pages between documents or change the page order: the page tree has to be
rebuilt. \texttt{rewritePdf} takes a plan --- for each output page, a source page index plus
optional rotation, crop box and a flatten step --- and produces a new file.

The invariant here is about \emph{inheritance}. In the source, a page may take
\texttt{MediaBox}, \texttt{CropBox}, \texttt{Resources}, \texttt{Rotate} and
\texttt{UserUnit} from an intermediate \texttt{/Pages} node. The rewrite produces a flat page
tree with one \texttt{/Pages} node, so each page must carry everything it inherited.
Omitting this does not fail: the file parses, and the pages are the wrong size, or sideways,
which is a far worse failure mode than an error.

Object handling follows from what the parser retained:

\begin{itemize}
  \item an object stored uncompressed is copied as its exact source text, with no
        re-serialisation, so anything the application does not model survives;
  \item an object inside an object stream is re-serialised from its parsed value, because its
        original bytes are inside a compressed container;
  \item a source page object is reused, and its dictionary rewritten, when that page appears
        exactly once in the plan; otherwise a fresh object is written;
  \item pages removed by the plan leave their objects unreferenced --- again legal orphans.
\end{itemize}

Because the plan can reference the same source page more than once, a rotation of a single
page is a one-page plan and a reorder is an identity plan with a permuted index list; the same
code path serves all of them.

\subsection{Merging two documents}
\label{sec:merge}

Merging is the one operation that copies between two object graphs. Every reference inside the
incoming document must be renumbered to its new home, or it will point at whatever the host
happens to keep at that number.

The algorithm is: build a map from each incoming object number to a fresh number above the
host's highest; for every incoming object except the page objects, its \texttt{/Pages} node
and its catalogue, re-serialise the value through a recursive reference rewrite, keeping
stream bytes and correcting \texttt{/Length} to the byte count; then build the incoming pages
on the host's page tree, with the same materialisation as \S\ref{sec:rewrite}.

Three details are easy to get wrong and were:

\begin{enumerate}
  \item \textbf{Stream bytes are not references.} A stream's data may contain any bytes,
        including ones that look like \texttt{N 0 R}. Only the dictionary is rewritten.
  \item \textbf{The incoming page tree is discarded, not merged.} Keeping the incoming
        \texttt{/Pages} node would leave a second tree with no parent, and readers differ on
        what to do with it.
  \item \textbf{Insertion is at a position, not an append.} Merging a cover sheet is the
        common case, so the interface takes a page index and clamps it to the document rather
        than rejecting it.
\end{enumerate}

\subsection{Flattening and summary documents}

\emph{Flattening} bakes annotations into page content: an extra content stream is appended
rather than merged into the existing one (a stream's bytes are immutable, and concatenating
inside one is a rewrite of bytes the writer promised not to touch), the page's
\texttt{/Annots} is dropped, and a Helvetica font is installed into a \emph{copy} of the
page's resource dictionary so the page's own font numbering is not disturbed.

\emph{Summarising} produces a new document containing the annotated pages. The obvious
implementation is a text listing of the comments, and it fails on the first document whose
comments are not in a Latin script: a text summary needs a font that covers whatever the
reader typed, and embedding one is a large amount of machinery for a summary. The
implementation instead renders each annotated page --- with its annotations already flattened
in --- at twice the page size and assembles the images into a new PDF. The result carries the
marks and the words inside them in any language, needs no font of its own, and costs nothing
that the printing path did not already provide.
